\section*{第 \ref{ch:g10:electron-shells} 章——电子层与元素周期表}

\begin{solution}{exo:g10:electron-shells:1}
C：$1\mathrm{s}^2\,2\mathrm{s}^2\,2\mathrm{p}^2$。N：
$1\mathrm{s}^2\,2\mathrm{s}^2\,2\mathrm{p}^3$。Mg：
$1\mathrm{s}^2\,2\mathrm{s}^2\,2\mathrm{p}^6\,3\mathrm{s}^2$。
\end{solution}

\begin{solution}{exo:g10:electron-shells:2}
碳 4，氮 5，镁 2。
\end{solution}

\begin{solution}{exo:g10:electron-shells:3}
最外\omterm{def:g10:electron-shells:configuration}{电子层} $n = 3$：第 3 周期。$2 + 3 = 5$ 个\omterm{def:g10:electron-shells:valence}{价电子}，属于
第 15 族（磷）。
\end{solution}

\begin{solution}{exo:g10:electron-shells:4}
\omterm{def:g10:electron-shells:family}{碱金属}（锂、钠、钾）；\omterm{def:g10:electron-shells:family}{卤素}（氟、
氯、溴、碘）；\omterm{def:g10:electron-shells:family}{稀有气体}（氦、氖、氩）。
\end{solution}

\begin{solution}{exo:g10:electron-shells:5}
\omterm{def:g7:atoms-and-molecules:atom}{原子}倾向于得失\omterm{def:g9:inside-the-atom:nucleus}{电子}，使最外\omterm{def:g10:electron-shells:configuration}{电子层}具有八个\omterm{def:g9:inside-the-atom:nucleus}{电子}，
与最近的\omterm{def:g10:electron-shells:family}{稀有气体}相同。
\end{solution}

\begin{solution}{exo:g10:electron-shells:6}
K，$2.8.8.1$，失去一个\omterm{def:g9:inside-the-atom:nucleus}{电子}：\ce{K+}，
$1\mathrm{s}^2\,2\mathrm{s}^2\,2\mathrm{p}^6\,3\mathrm{s}^2\,3\mathrm{p}^6$
（氩）。F，$2.7$，得到一个\omterm{def:g9:inside-the-atom:nucleus}{电子}：\ce{F-}，
$1\mathrm{s}^2\,2\mathrm{s}^2\,2\mathrm{p}^6$（氖）。
\end{solution}

\begin{solution}{exo:g10:electron-shells:7}
$1\mathrm{s}^2\,2\mathrm{s}^2\,2\mathrm{p}^6\,3\mathrm{s}^2$：$Z = 12$，
为镁。
\end{solution}

\begin{solution}{exo:g10:electron-shells:8}
\ce{Mg^{2+}} 与 \ce{Cl-}：\ce{MgCl2}。\ce{Al^{3+}} 与 \ce{O^{2-}}：
\ce{Al2O3}。
\end{solution}

\begin{solution}{exo:g10:electron-shells:9}
硫（6 个\omterm{def:g10:electron-shells:valence}{价电子}，$2.8.6$），同属第 16 族。
\end{solution}

\begin{solution}{exo:g10:electron-shells:10}
其最外\omterm{def:g10:electron-shells:configuration}{电子层}已满（八个\omterm{def:g9:inside-the-atom:nucleus}{电子}），因此
无论得失\omterm{def:g9:inside-the-atom:nucleus}{电子}都没有好处。
\end{solution}

\begin{solution}{exo:g10:electron-shells:11}
氩有 18 个\omterm{def:g9:inside-the-atom:nucleus}{电子}；\ce{X^{2+}} 失去了 2 个，因此 X 有 20 个\omterm{def:g9:inside-the-atom:nucleus}{电子}：钙。
\end{solution}

\begin{solution}{exo:g10:electron-shells:12}
氦，$1\mathrm{s}^2$，最外\omterm{def:g10:electron-shells:configuration}{电子层}已满（第 1 \omterm{def:g10:electron-shells:configuration}{电子层}只能容纳两个
\omterm{def:g9:inside-the-atom:nucleus}{电子}），与其他\omterm{def:g10:electron-shells:family}{稀有气体}相同；它的性质不同于第 2 族
\omterm{def:g9:periodic-table-first-look:metal}{金属}。
\end{solution}

\begin{solution}{exo:g10:electron-shells:13}
其\omterm{def:g10:electron-shells:valence}{价电子}位于第 4 \omterm{def:g10:electron-shells:configuration}{电子层}，比钠的
\omterm{def:g10:electron-shells:valence}{价电子}（位于第 3 \omterm{def:g10:electron-shells:configuration}{电子层}）离\omterm{def:g9:inside-the-atom:nucleus}{原子核}更远，束缚较弱，因此更容易失去。
\end{solution}

\begin{solution}{exo:g10:electron-shells:14}
这四者均有 18 个\omterm{def:g9:inside-the-atom:nucleus}{电子}，\omterm{def:g10:electron-shells:configuration}{电子排布}与氩相同：
\[ 1\mathrm{s}^2\,2\mathrm{s}^2\,2\mathrm{p}^6\,3\mathrm{s}^2\,3\mathrm{p}^6. \]
\end{solution}

\begin{solution}{exo:g10:electron-shells:15}
它得到 3 个\omterm{def:g9:inside-the-atom:nucleus}{电子}以达到\omterm{def:g10:electron-shells:octet}{八隅规则}：它有 5 个\omterm{def:g10:electron-shells:valence}{价电子}，
属第 15 族。位于第 3 周期，即磷：
\[ 1\mathrm{s}^2\,2\mathrm{s}^2\,2\mathrm{p}^6\,3\mathrm{s}^2\,3\mathrm{p}^3. \]
\end{solution}

\begin{solution}{pb:g10:electron-shells:1}
\textbf{1.} 铝：13 个\omterm{def:g9:inside-the-atom:proton}{质子}，13 个\omterm{def:g9:inside-the-atom:nucleus}{电子}。氧：8 个\omterm{def:g9:inside-the-atom:proton}{质子}，8
个\omterm{def:g9:inside-the-atom:nucleus}{电子}。

\textbf{2.} Al：$1\mathrm{s}^2\,2\mathrm{s}^2\,2\mathrm{p}^6\,3\mathrm{s}^2\,3\mathrm{p}^1$。
O：$1\mathrm{s}^2\,2\mathrm{s}^2\,2\mathrm{p}^4$。

\textbf{3.} 铝：3 个\omterm{def:g10:electron-shells:valence}{价电子}，第 3 周期，第 13 族。氧：6 个\omterm{def:g10:electron-shells:valence}{价电子}，第 2 周期，
第 16 族。

\textbf{4.} 铝是\omterm{def:g9:periodic-table-first-look:metal}{金属}，氧是\omterm{def:g9:periodic-table-first-look:metal}{非金属}。

\textbf{5.} \ce{Al^{3+}}：$1\mathrm{s}^2\,2\mathrm{s}^2\,2\mathrm{p}^6$。

\textbf{6.} \ce{O^{2-}}：$1\mathrm{s}^2\,2\mathrm{s}^2\,2\mathrm{p}^6$。

\textbf{7.} 氖。

\textbf{8.} 铝失去 3 个\omterm{def:g9:inside-the-atom:nucleus}{电子}，氧得到 2 个\omterm{def:g9:inside-the-atom:nucleus}{电子}。

\textbf{9.} $2 \times (+3) + 3 \times (-2) = 0$：\ce{Al2O3}。

\textbf{10.} 失去：$2 \times 3 = 6$。得到：$3 \times 2 = 6$。

\textbf{11.} \omterm{def:g9:inside-the-atom:nucleus}{电子}既不会产生也不会消失：铝
失去的\omterm{def:g9:inside-the-atom:nucleus}{电子}恰好等于氧得到的\omterm{def:g9:inside-the-atom:nucleus}{电子}，化合物
呈电中性。

\textbf{12.} 它所带电荷相同，为 $+3$，晶体中的
电荷仍然平衡。

\textbf{13.} $2 \times 27 + 3 \times 16 = 102$ 个\omterm{def:g9:inside-the-atom:proton}{核子}。

\textbf{14.} $102 \times 1.67 \times 10^{-27} = \qty{1.70e-25}{kg}$。

\textbf{15.} $10^{-3} / 1.70 \times 10^{-25} \approx 5.87 \times 10^{21}$
化学式单元。

\textbf{16.} $2 \times 5.87 \times 10^{21} = 1.17 \times 10^{22}$
\ce{Al^{3+}} \omterm{def:g9:ions:ion}{离子}；$3 \times 5.87 \times 10^{21} = 1.76 \times 10^{22}$
\ce{O^{2-}} \omterm{def:g9:ions:ion}{离子}。

\textbf{17.} \omterm{def:g9:inside-the-atom:nucleus}{电子}约比\omterm{def:g9:inside-the-atom:proton}{核子}轻 1836 倍；
一个化学式单元的 50 个\omterm{def:g9:inside-the-atom:nucleus}{电子}，
总质量不足一个\omterm{def:g9:inside-the-atom:proton}{核子}的百分之三。

\textbf{18.} $6 \times 5.87 \times 10^{21} \approx 3.5 \times 10^{22}$
个\omterm{def:g9:inside-the-atom:nucleus}{电子}转移，对应 \qty{1}{g} 刚玉。
\end{solution}
